
\documentclass[12pt]{article}
\usepackage[utf8]{inputenc}
\usepackage{amsmath}
\usepackage{graphicx}
\usepackage{listings}
\usepackage{hyperref}
\usepackage{geometry}
\usepackage{booktabs}
\geometry{margin=1in}

\title{Lab 2: Sequence-to-Sequence Neural Machine Translation}
\author{}
\date{}

\begin{document}
\maketitle

\section{Objective}
The goal of this lab is to implement a basic sequence-to-sequence (Seq2Seq)
model for neural machine translation.

After completing this lab, students should be able to:
\begin{itemize}
    \item Explain the role of the encoder and decoder.
    \item Implement a GRU-based Seq2Seq model.
    \item Train the model with teacher forcing.
    \item Perform greedy decoding at inference time.
    \item Evaluate translations using a simple exact-match accuracy or BLEU score.
\end{itemize}

\section{Background}
A Seq2Seq model maps an input sentence
\[
x = (x_1, x_2, \dots, x_T)
\]
to an output sentence
\[
y = (y_1, y_2, \dots, y_{T'}).
\]

In the basic encoder--decoder formulation:
\begin{itemize}
    \item The \textbf{encoder} reads the source sentence and produces a hidden representation.
    \item The \textbf{decoder} generates the target sentence token by token.
\end{itemize}

For this lab, we use GRU-based encoder and decoder modules.

\section{Dataset}
The dataset is a small English--Vietnamese parallel corpus stored in:
\begin{itemize}
    \item \texttt{data/train.en}
    \item \texttt{data/train.vi}
    \item \texttt{data/dev.en}
    \item \texttt{data/dev.vi}
\end{itemize}

Students may reuse the vocabulary-building utilities from Lab 1.

\section{Architecture}
The model is defined as:
\begin{enumerate}
    \item Embed source tokens.
    \item Encode the source sequence using a GRU.
    \item Use the final encoder hidden state to initialize the decoder.
    \item Generate target tokens step by step using the decoder GRU.
    \item Apply a linear layer to obtain token logits.
\end{enumerate}

\section{Task 1: Data Preparation}
Implement or reuse the following components:
\begin{itemize}
    \item tokenization
    \item vocabulary construction
    \item sentence encoding with \texttt{<BOS>} and \texttt{<EOS>}
    \item mini-batch padding
\end{itemize}

\subsection*{Question}
Why do we need \texttt{<BOS>} and \texttt{<EOS>} tokens in sequence generation?

\section{Task 2: Encoder}
Implement a GRU encoder:

\begin{lstlisting}[language=Python]
class Encoder(nn.Module):
    def __init__(self, vocab_size, emb_dim, hidden_dim):
        ...
    def forward(self, x):
        ...
\end{lstlisting}

The encoder should return:
\begin{itemize}
    \item all hidden states (optional for later labs)
    \item the final hidden state
\end{itemize}

\section{Task 3: Decoder}
Implement a GRU decoder:

\begin{lstlisting}[language=Python]
class Decoder(nn.Module):
    def __init__(self, vocab_size, emb_dim, hidden_dim):
        ...
    def forward(self, token, hidden):
        ...
\end{lstlisting}

At each decoding step, the decoder should:
\begin{enumerate}
    \item embed the previous token,
    \item update the hidden state,
    \item compute logits over the target vocabulary.
\end{enumerate}

\subsection*{Question}
What probability distribution is modeled by the decoder at time step $t$?

\section{Task 4: Seq2Seq Training}
Implement the Seq2Seq forward pass with teacher forcing.

Training objective:
\[
L = - \sum_{t=1}^{T'} \log p(y_t \mid y_{<t}, x)
\]

Use cross-entropy loss and train the model for several epochs.

\subsection*{Question}
What is teacher forcing? Why is it useful during training?

\section{Task 5: Greedy Decoding}
At inference time, generate the translation token by token using:
\[
\hat{y}_t = \arg\max_y p(y \mid \hat{y}_{<t}, x)
\]

Stop generation when:
\begin{itemize}
    \item \texttt{<EOS>} is produced, or
    \item a maximum target length is reached.
\end{itemize}

\section{Task 6: Evaluation}
Evaluate the model on the development set.
You may report:
\begin{itemize}
    \item exact-match accuracy on the toy dataset, or
    \item BLEU score if you choose to use NLTK.
\end{itemize}

\subsection*{Suggested output}
For a few development examples, print:
\begin{itemize}
    \item source sentence
    \item reference translation
    \item model prediction
\end{itemize}

\section{Submission}
Students must submit:
\begin{itemize}
    \item \texttt{lab2_seq2seq.ipynb}
    \item a short PDF with answers to conceptual questions
\end{itemize}

\section{Grading}
\begin{center}
\begin{tabular}{lc}
\toprule
Task & Points \\
\midrule
Data preparation & 2 \\
Encoder implementation & 2 \\
Decoder implementation & 2 \\
Seq2Seq forward + training & 2 \\
Greedy decoding + evaluation & 2 \\
\midrule
Total & 10 \\
\bottomrule
\end{tabular}
\end{center}

\section{Expected Behavior}
Because the dataset is very small, the model may overfit quickly.
A correct implementation should show:
\begin{itemize}
    \item decreasing training loss,
    \item reasonable translations for memorized sentence patterns,
    \item better training behavior when teacher forcing is used.
\end{itemize}

\end{document}
